% ============================================================
% v2 — Capitolo 6: Conclusioni
% Blocco IV di IV. Budget: ~2 pagine.
% ============================================================
\chapter{Conclusioni}
\label{ch:conclusioni}

Questa tesi ha chiesto se ordinare i task di addestramento dal facile al
difficile produca un comportamento di guida multi-agente misurabilmente
diverso rispetto al presentare tutte le difficoltà insieme, a parità di
budget, in un ambiente CARLA a mondo condiviso addestrato con MAPPO. La
risposta è sì --- ma non lungo l'asse su cui i curricula vengono di
solito proposti.

\paragraph{L'ordinamento non ha accelerato l'apprendimento.}
Tutte e quattro le celle della campagna hanno raggiunto le stesse soglie
iniziali di competenza agli stessi conteggi di passi, con uno scarto
sotto l'8\% (\cref{sec:efficienza}), e questo pur avendole superate su
mix di difficoltà completamente diversi. La spiegazione più plausibile è
che il reward shaping denso di questa piattaforma fornisca già segnale
di apprendimento anche sui task difficili, rendendo superfluo l'avvio
graduale: una condizione di scope su H1, non una confutazione del
curriculum learning in generale.

\paragraph{L'ordinamento ha selezionato ciò che la policy è diventata.}
Sotto il critic MLP il curriculum ha prodotto un guidatore cauto e
incline allo stallo e il batch uno veloce e incline alle collisioni; sul
criterio di successo rigoroso lo scambio si è risolto a favore del
curriculum di $+18.5$\,pp in valutazione, coerentemente in tutti e
quattro gli scenari (\cref{sec:prestazione}). L'ordinamento dei task
agisce cioè da \emph{selettore di equilibri} sul trade-off
velocità--sicurezza, più che da acceleratore.

\paragraph{L'effetto è condizionato dalla stima del valore.}
Sotto un critic GNN l'effetto di regime in-domain è svanito
($-0.1$\,pp) mentre l'architettura ha recuperato la cella batch e
rimodellato la cella curriculum: una forte interazione
architettura$\times$regime (\cref{sec:asse-arch}). Il regime di
addestramento non è quindi una leva indipendente dal resto dello stack,
e riportarne l'effetto senza dichiarare l'architettura del critic è
riportarne solo metà.

\paragraph{Il beneficio portabile è la generalizzazione.}
L'unica componente del vantaggio del curriculum replicata in entrambe le
architetture è il margine su una città mai vista in addestramento
($+19.0$ e $+10.0$\,pp su Town05), ottenuta contro l'intuizione
copertura-implica-robustezza che motivava H2
(\cref{sec:generalizzazione}). È il risultato che questa tesi difende
con più forza, perché è l'unico osservato due volte in modo
indipendente.

\paragraph{Specificità di classe e limiti della piattaforma.}
I pedoni sono risultati in larga parte insensibili al regime, vincolati
dalla fattibilità dei percorsi più che dalla qualità della policy
(\cref{sec:pedoni}); e ogni cella ha mostrato erosione tardiva entro la
run a livello di task fissato, indipendente dal regime
(\cref{sec:dinamica}). Le scelte di regime calibrate sulla classe forte
non vanno assunte trasferibili alla debole.

\bigskip

Oltre ai risultati, la tesi contribuisce l'apparato che li ha resi
misurabili: un framework di confronto a parità di budget con
configurazioni appaiate per seed e byte-identiche; un protocollo di
sviluppo guidato da gate con un registro esplicito di modifiche
promosse, rigettate e consapevolmente mantenute; un resoconto
documentato di come i difetti di validità dell'ambiente possano
mascherarsi da effetti di policy, con le loro firme di rilevamento; e un
metro di varianza da repliche a seed identico contro cui ogni
affermazione viene letta. In un contesto in cui un singolo bug del
simulatore ha spostato la metrica principale di cinquanta punti
percentuali senza che nulla cambiasse nella policy, questa disciplina di
misura --- e non un singolo numero --- è ciò che rende il confronto
significativo.

Resta il limite che la tesi dichiara apertamente e che il
\cref{ch:evoluzioni} mette in cima all'elenco: con una run per cella,
ogni contrasto è un caso di studio appaiato letto contro una banda di
rumore, non la stima di una distribuzione. La replicazione sui seed
della matrice $2\times2$ è il passo che trasformerebbe le conclusioni
qui riportate da direzioni ben corroborate in effetti quantificati.
